\begin{abstract}
A marker vanishes from a diseased tissue: did the cells making it leave, or did
the survivors turn it down? Single-cell data promises to separate the two, and we
built that measurement for colorectal cancer behind a pre-registered control
panel, a confound diagnostic, a simulation-based validation sweep and an
abstention rule. Every guard passed. \textbf{A check unable to fail is worse than
no check, because it turns an absence of evidence into a green light}, and four
of ours could not fail.

Two of the four generalise past our data and are provable rather than observed.
A threshold on the correlation between sequencing depth and a \emph{binary} cell
label of prevalence $p$ cannot fire at all once the threshold exceeds
$\sqrt{3p(1-p)}$ --- so on rare cell types, the populations worth resolving, the
diagnostic reports clean whatever the data does. And a recovery curve cannot see
an estimator that is a deterministic function of the statistics its simulator
drew: the estimator cancels exactly, and the curve measures the generator. Each
is one line to check and neither is specific to biology, though what forces us
into both is: a leakage constraint bars a cell type's positive markers from
labelling it when those same markers are the quantity being averaged, so the
label is defined by \emph{absence}, and absence reads dropout.

We give the checks we now run, abstention as a first-class output, and the
result that abstention's own calibration does not return the cutpoint we
committed to.
\end{abstract}
